\section{Related Work}
\label{sec:related}

\subsection{Deep Learning for Plant Disease Detection}

Interest in deep learning for plant disease detection has grown
substantially over the past decade, progressing from laboratory-scale
binary classifiers to multi-class systems tested under real field
conditions.

\subsubsection*{Foundational CNN-based classification.}
Mohanty et al.~\cite{mohanty2016} were among the first to demonstrate
that CNNs trained on the PlantVillage dataset could exceed 99\%
classification accuracy, establishing transfer learning as the
dominant strategy in the field.
Geetharamani and Pandian~\cite{geetharamani2019} developed a
nine-layer deep CNN that achieved precision and recall of 99.47\% and
99.89\% respectively, though their evaluation was confined to images
with clean, uniform backgrounds.
Ferentinos~\cite{ferentinos2018} extended this line of work to
25 disease types and reported accuracy above 99.5\% under controlled
studio conditions, but documented a marked performance decline on
images from actual fields, a gap that remains relevant to any
real-world deployment.

\subsubsection*{Object detection for disease localisation.}
Classification identifies whether disease is present but cannot
pinpoint where. For precision treatment, spatial localisation matters.
Nawaz et al.~\cite{nawaz2023} developed a VGG-19-based Faster RCNN
system achieving mAP = 0.87 on PlantVillage, outperforming several
competing baselines (Fast-RCNN/VGG-16: 0.85; YOLOv3/DarkNet-53:
0.83; SSD/ResNet-101: 0.81).
The Regional Proposal Network inside Faster RCNN uses shape, colour,
and structural cues to localise diseased tissue, which directly
motivated our adoption of this architecture in Stage~2.

\subsubsection*{EfficientNet for agricultural image analysis.}
The EfficientNet family~\cite{efficientnet} applies compound scaling
to width, depth and resolution simultaneously, producing strong
accuracy-per-parameter in agricultural classification tasks.
Kadry et al.~\cite{kadry2025} applied variants B0 to B4 to farm insect
classification, with fused features reaching 98.3\% accuracy.
Das et al.~\cite{das2024} reviewed the broader CNN literature and
found that modern families such as ResNet, EfficientNet and ConvNeXt
consistently outperform hand-crafted feature approaches under variable
colour and illumination.

\subsubsection*{Explainability.}
Grad-CAM~\cite{selvaraju2017} produces coarse localisation maps from
the gradients flowing into the last convolutional layer, and SHAP~\cite{lundberg2017}
attributes a model's output to its input features using Shapley
values. In plant pathology, such tools matter because a classifier
that reaches high accuracy by attending to background, leaf edges or
acquisition artefacts will fail in the field; Argüeso et
al.~\cite{arguesoa2020}, for example, studied exactly this gap between
laboratory and field imagery. We use Grad-CAM to inspect the
classifier's evidence and SHAP to audit the yield model.

\subsubsection*{Transformer-based and hybrid approaches.}
Vision Transformers~\cite{dosovitskiy2021} and hierarchical variants
such as Swin~\cite{liu2021swin} have become strong alternatives to
CNNs on large benchmarks, typically relying on large-scale
pre-training. Our benchmark is restricted to the CNN architectures available in the
TensorFlow 2.10 Keras applications library, trained with a shared head,
optimiser and augmentation, with the exceptions described in
Section~\ref{sec:training}. This excludes transformers and also more recent
CNNs such as ConvNeXt V2, which are natural extensions left to future work. Hybrid frameworks
that combine learned and process-based components for yield
prediction are discussed in Section~\ref{sec:integrated}.

\subsubsection*{Remaining gaps.}
To our knowledge, no prior study has benchmarked all major CNN
families on a single diverse crop disease dataset under a unified
experimental protocol.
More critically, the connection between a visual disease classifier
and a downstream yield prediction model has not been formalised or
tested as an integrated pipeline.

\subsection{Automated Annotation}

Generating bounding box annotations at scale is one of the primary
bottlenecks in building plant pathology detection datasets.
Nawaz et al.~\cite{nawaz2023} relied on LabelImg for manual
annotation, a process requiring domain expertise that is difficult to
scale beyond a few hundred images.
Grounding DINO~\cite{groundingdino} provides a practical alternative,
localising objects described only in plain text with no labelled
training examples required.
Liu et al.~\cite{liu2019pestnet} and Dewari et
al.~\cite{dewari2022} both identified annotation volume as the
principal limiting factor in agricultural detection, motivating our
use of Grounding DINO to expand the annotated set from 280 to 2,256
images without any manual effort.

\subsection{Crop Yield Prediction}

Ensemble regressors such as XGBoost and Random Forest have
consistently outperformed linear methods on crop yield prediction
tasks involving complex climate-yield
interactions~\cite{ml_yield_review}.
A recurring finding in the literature is that incorporating disease
incidence data into yield models can reduce prediction error by as
much as 18\%~\cite{disease_yield_integration}.
To our knowledge, prior work has rarely combined image-derived disease
information with yield prediction in a single framework; we do not claim a
validated image-to-yield system, since our yield records are not paired
with the leaf images.

\subsection{Integrated Approaches}
\label{sec:integrated}

Several groups have combined process-based crop simulation with
machine learning to improve yield forecasting under climate
change~\cite{hybrid_yield}.
Our work takes this integration further by linking the outputs of a
vision-based disease detector directly to the yield prediction step,
forming a three-stage image-to-severity-to-yield framework, although the
disease imagery and yield records are not currently paired by field, season
or location.

